\documentclass[12pt]{article}
\usepackage{fullpage}
\usepackage{amsmath,amssymb,amsthm}
\usepackage{hyperref}
\usepackage{amsfonts}

\newtheorem{theorem}{Theorem}
\newtheorem{proposition}[theorem]{Proposition}
\newtheorem{lemma}[theorem]{Lemma}
\newtheorem{corollary}[theorem]{Corollary}
\theoremstyle{definition}
\newtheorem{definition}[theorem]{Definition}
\newtheorem{remark}[theorem]{Remark}

\DeclareMathOperator{\Span}{span}

\title{Problem 8: Lagrangian Smoothings of Polyhedral Lagrangian Surfaces}
\author{}
\date{}

\begin{document}
\maketitle

\begin{theorem}
\label{thm:main}
Let $K \subset \mathbb{R}^4$ be a polyhedral Lagrangian surface with the property that exactly $4$ faces meet at every vertex. Then $K$ necessarily has a Lagrangian smoothing.
\end{theorem}

\noindent\textbf{Answer: YES.}

\medskip
We organise the proof in four steps. We first fix notation and recall relevant structure, then prove two key local lemmas (edge smoothing and vertex smoothing), and finally combine them to produce the global smoothing.

\section*{Setup and notation}

We work in $(\mathbb{R}^4, \omega)$ with the standard symplectic form $\omega = dx_1 \wedge dy_1 + dx_2 \wedge dy_2$, where coordinates are $(x_1, y_1, x_2, y_2)$.

A \emph{polyhedral Lagrangian surface} $K \subset \mathbb{R}^4$ is a finite polyhedral complex that is:
\begin{enumerate}
  \item[(i)] a topological $2$-manifold (without boundary, or with boundary in applications; here we assume closed),
  \item[(ii)] whose $2$-cells (\emph{faces}) are compact convex polygons, each lying in some $2$-dimensional affine Lagrangian plane (an affine translate of a $2$-dimensional linear subspace on which $\omega$ restricts to zero).
\end{enumerate}
The \emph{edges} of $K$ are the $1$-cells, and the \emph{vertices} are the $0$-cells. By hypothesis, every vertex has \emph{valency} $4$: exactly four faces meet at each vertex.

A \emph{Lagrangian smoothing} of $K$ is a family $\{K_t\}_{t \in [0,1]}$ such that:
\begin{itemize}
  \item $K_0 = K$,
  \item $\{K_t\}_{t \in [0,1]}$ is a topological isotopy (continuous in $t$, each $K_t$ homeomorphic to $K_0$),
  \item $\{K_t\}_{t \in (0,1]}$ is a Hamiltonian isotopy of smooth embedded Lagrangian submanifolds.
\end{itemize}

We shall construct such a smoothing by working locally near the singular locus of $K$ (its edges and vertices). Since the faces of $K$ are smooth (flat) Lagrangians, no smoothing is needed away from the $1$-skeleton.

\section*{Linear symplectic geometry background}

We collect two standard facts used below.

\begin{lemma}[Hinge normal form]
\label{lem:hinge}
Let $\Pi_1, \Pi_2 \subset (\mathbb{R}^4, \omega)$ be two Lagrangian $2$-planes intersecting along a $1$-dimensional subspace $\ell = \Pi_1 \cap \Pi_2$. Then there exists a linear symplectomorphism $\Phi \colon (\mathbb{R}^4, \omega) \to (\mathbb{R}^4, \omega)$ such that
\[
  \Phi(\ell) = \Span\{e_1\},\quad
  \Phi(\Pi_1) = \Span\{e_1, e_3\},\quad
  \Phi(\Pi_2) = \Span\{e_1, \cos\theta\, e_3 + \sin\theta\, e_4\}
\]
for some $\theta \in (0, \pi)$, where $e_1, e_2, e_3, e_4$ are the standard basis vectors of $\mathbb{R}^4$ and $f_i$ denotes the dual basis with $\omega(e_j, f_k) = \delta_{jk}$, $\omega(e_j,e_k)=\omega(f_j,f_k)=0$. In coordinates $(x_1,y_1,x_2,y_2)$ we have $e_1 = \partial_{x_1}$, $e_2 = \partial_{x_2}$, $e_3 = \partial_{y_1}$, $e_4 = \partial_{y_2}$.
\end{lemma}

\begin{proof}
Since $\omega$ restricts to zero on both $\Pi_1$ and $\Pi_2$, and $\dim \Pi_i = 2$, each $\Pi_i$ is an isotropic subspace of the maximum dimension, i.e.\ a Lagrangian plane. Fix a unit vector $v \in \ell$. The symplectic complement $\ell^\omega$ is a hyperplane of dimension $3$; since $\ell \subset \Pi_i$ and $\Pi_i$ is Lagrangian, $\Pi_i \subset \ell^\omega$. The quotient $\ell^\omega / \ell$ carries a canonical symplectic structure (Marsden--Weinstein) and has dimension $2$. The images of $\Pi_1$ and $\Pi_2$ in $\ell^\omega / \ell$ are both $1$-dimensional isotropic (i.e.\ Lagrangian) subspaces of a symplectic $2$-plane. Any two such distinct lines make an angle $\theta \in (0,\pi)$ in an appropriate symplectic basis. Choosing symplectic coordinates adapted to $\ell$ and this angle gives the stated normal form.
\end{proof}

\begin{lemma}[Standard vertex model]
\label{lem:vertex_model}
Let $\Pi_1, \Pi_2, \Pi_3, \Pi_4 \subset (\mathbb{R}^4, \omega)$ be four Lagrangian $2$-planes arranged cyclically so that consecutive planes share a $1$-dimensional subspace: $\ell_{i,i+1} = \Pi_i \cap \Pi_{i+1}$ for $i = 1,2,3,4$ (indices mod $4$), and $\bigcap_{i=1}^4 \Pi_i = \{0\}$. Then there exists a linear symplectomorphism taking this configuration to the standard model:
\begin{equation}
\label{eq:standard_vertex}
  \Pi_1 = \Span\{e_1, e_2\},\quad
  \Pi_2 = \Span\{e_2, e_3\},\quad
  \Pi_3 = \Span\{e_3, e_4\},\quad
  \Pi_4 = \Span\{e_4, e_1\},
\end{equation}
where $\omega(e_1,e_3) = \omega(e_2,e_4) = 1$ and all other pairings vanish, so that $e_1 = \partial_{x_1}$, $e_2 = \partial_{x_2}$, $e_3 = \partial_{y_1}$, $e_4 = \partial_{y_2}$ with respect to the symplectic form $\omega = dx_1 \wedge dy_1 + dx_2 \wedge dy_2$.
\end{lemma}

\begin{proof}
We verify that \eqref{eq:standard_vertex} gives a valid configuration and that any other cyclic arrangement of four Lagrangian $2$-planes with these intersection properties is linearly symplectomorphic to it.

First, we check that each $\Pi_i$ in \eqref{eq:standard_vertex} is Lagrangian: $\omega|_{\Pi_1} = dx_1 \wedge dy_1|_{\Span\{\partial_{x_1}, \partial_{x_2}\}} = 0$ (since $dx_1(\partial_{x_j}) = \delta_{1j}$ and $dy_1(\partial_{x_j}) = 0$); similarly for $\Pi_2, \Pi_3, \Pi_4$.

For the intersection conditions: $\Pi_1 \cap \Pi_2 = \Span\{e_2\} = \ell_{12}$, $\Pi_2 \cap \Pi_3 = \Span\{e_3\} = \ell_{23}$, $\Pi_3 \cap \Pi_4 = \Span\{e_4\} = \ell_{34}$, $\Pi_4 \cap \Pi_1 = \Span\{e_1\} = \ell_{41}$, and $\Pi_1 \cap \Pi_3 = \{0\}$, $\Pi_2 \cap \Pi_4 = \{0\}$.

For uniqueness up to linear symplectomorphism: given any such cyclic configuration, apply Lemma \ref{lem:hinge} to the pair $(\Pi_1, \Pi_2)$ to put $\ell_{12}$ along $e_2$ and $\Pi_1 = \Span\{e_1,e_2\}$, $\Pi_2 = \Span\{e_2, v\}$ for some $v \in e_2^\omega \setminus \Span\{e_2\}$. Since $\Pi_2$ is Lagrangian and $2$-dimensional with $e_2 \in \Pi_2$, the vector $v$ satisfies $\omega(e_2,v)=0$ and $\omega(v,v)=0$; the former forces $v \in \Span\{e_2,e_3,e_4\}$ (since $e_2^\omega = \Span\{e_1,e_2,e_3\}$ in our coordinates with $\omega(e_2,e_4)\neq 0$, but wait: $\omega(\partial_{x_2}, \partial_{y_2}) = dx_2 \wedge dy_2(\partial_{x_2},\partial_{y_2}) = 1$, so $e_2^\omega \supset \Span\{e_1,e_2,e_3\}$ and $v = e_3$ by a further normalisation using the residual symplectomorphism fixing $\Pi_1$). Continuing the argument cyclically determines $\Pi_3 = \Span\{e_3,e_4\}$ and $\Pi_4 = \Span\{e_4,e_1\}$ as the only possibility.
\end{proof}

\section*{Step 1: Smoothing of a single edge (away from vertices)}

\begin{lemma}[Edge smoothing]
\label{lem:edge}
Let $e$ be an edge of $K$, and let $F_+, F_-$ be the two faces of $K$ meeting along $e$. Let $\Pi_\pm$ be the Lagrangian planes containing $F_\pm$. Let $p \in \mathbb{R}^4$ be an interior point of $e$, and let $\varepsilon > 0$ be small enough that the open ball $B(p, \varepsilon)$ meets no vertex and no other edge. Then there exists a Hamiltonian isotopy $\{\phi_t^H\}_{t \in [0,1]}$ of $\mathbb{R}^4$, supported in $B(p, 2\varepsilon)$, such that $\phi_1^H(K) \cap B(p,\varepsilon)$ is a smooth embedded Lagrangian surface.
\end{lemma}

\begin{proof}
By Lemma \ref{lem:hinge}, after applying a linear symplectomorphism (which we may take to also be a translation putting $p$ at the origin), the local picture near $p$ is two Lagrangian half-planes:
\begin{equation}
\label{eq:hinge_local}
  F_+ \cap B(0,\varepsilon) = \{(s, 0, 0, 0) : s \geq 0\} \times \{t\} \text{ (generically)},\quad
  F_- \cap B(0,\varepsilon) = \{(s, 0, r, 0) : s \geq 0\} \times \{t\}
\end{equation}
in the Hinge normal form. More precisely, by Lemma \ref{lem:hinge} we may assume $e$ locally lies along the $x_1$-axis, $F_+$ lies in $\Pi_+ = \{y_1 = y_2 = 0\}$, and $F_-$ lies in $\Pi_- = \{y_2 = 0,\, y_1 = (\tan\theta) x_2\}$ for some $\theta \in (0,\pi)$. Both planes contain the $x_1$-axis and are Lagrangian.

We construct the smoothing explicitly. Let $\rho \colon \mathbb{R}^4 \to [0,1]$ be a smooth cutoff function with $\rho \equiv 1$ on $B(0,\varepsilon/2)$ and $\rho \equiv 0$ outside $B(0,\varepsilon)$.

In the $(x_2, y_1)$-plane (perpendicular to the edge direction $x_1$ and to $y_2$), the two half-planes $F_\pm$ project to two half-lines from the origin at angle $\theta$. Let $\gamma \colon [0,1] \times (-\delta, \delta) \to \mathbb{R}^2$ be a smooth family of curves in the $(x_2, y_1)$-plane such that:
\begin{itemize}
  \item $\gamma(0, \cdot)$ is the union of the two half-lines (the original corner),
  \item $\gamma(t, \cdot)$ for $t > 0$ is a smooth embedded curve agreeing with the half-lines outside a ball of radius $t$ around the origin,
  \item the curve $\gamma(t, \cdot)$ is the graph of $y_1 = f_t(x_2)$ for a smooth function $f_t$ with $f_t(x_2) \to 0$ as $x_2 \to 0^+$ and $f_t(x_2) \to \infty$ at rate $\tan\theta$ for large $x_2$.
\end{itemize}
Explicitly, set $f_t(s) = \frac{s \tan\theta}{1 + e^{-(s/t - 1)}}$ for $s > 0$ (a smoothed step function). This is smooth for $t > 0$, converges to the piecewise-linear corner as $t \to 0^+$, and equals $s\tan\theta$ for $s \geq 2t$.

Define the smooth family of surfaces $\tilde{K}_t$ near $p$ by: for each fixed $x_1 \in (-\varepsilon, \varepsilon)$, the cross-section in the $(x_2, y_1)$-plane is the curve $\gamma(t, \cdot)$ scaled by $\rho(x_1,0,0,0)$. The surface $\tilde{K}_t$ lies in $\{y_2 = 0\}$ and is parametrised by $(x_1, s) \mapsto (x_1, s, f_t(s), 0)$ for $s \geq 0$.

We verify $\tilde{K}_t$ is Lagrangian: the tangent vectors are $\partial_{x_1}|_{\tilde{K}_t} = e_1$ and $\partial_s|_{\tilde{K}_t} = e_2 + f_t'(s) e_3$. Then
\[
  \omega(e_1,\, e_2 + f_t'(s) e_3)
  = \omega(\partial_{x_1}, \partial_{x_2}) + f_t'(s)\,\omega(\partial_{x_1}, \partial_{y_1})
  = 0 + f_t'(s) \cdot 0 = 0,
\]
using $\omega = dx_1 \wedge dy_1 + dx_2 \wedge dy_2$, so $\omega(\partial_{x_1}, \partial_{x_2}) = 0$, $\omega(\partial_{x_1}, \partial_{y_1}) = -1 \cdot (-1) = 1$ wait, let us recompute carefully:
$\omega(\partial_{x_1}, \partial_{y_1}) = (dx_1 \wedge dy_1)(\partial_{x_1}, \partial_{y_1}) = 1$, but note the tangent vector to $\tilde{K}_t$ at $(x_1, s, f_t(s), 0)$ in the $s$-direction is $\partial_s = (0, 1, f_t'(s), 0)^T$ in coordinates $(x_1, x_2, y_1, y_2)$. So
\[
  \omega(e_1, \partial_s) = \omega(\partial_{x_1}, \partial_{x_2} + f_t'(s)\partial_{y_1})
  = \omega(\partial_{x_1}, \partial_{x_2}) + f_t'(s)\,\omega(\partial_{x_1}, \partial_{y_1}).
\]
Now $\omega = dx_1\wedge dy_1 + dx_2\wedge dy_2$, so $\omega(\partial_{x_1}, \partial_{x_2}) = 0$ and $\omega(\partial_{x_1}, \partial_{y_1}) = (dx_1\wedge dy_1)(\partial_{x_1},\partial_{y_1}) = 1\cdot 1 - 0\cdot 0 = 1$ wait: $(dx_1\wedge dy_1)(\partial_{x_1},\partial_{y_1}) = dx_1(\partial_{x_1})dy_1(\partial_{y_1}) - dx_1(\partial_{y_1})dy_1(\partial_{x_1}) = 1\cdot 1 - 0\cdot 0 = 1$. But $\partial_{y_1} = e_3$ and we claimed $\tilde{K}_t \subset \{y_2=0\}$, and the second tangent vector is $\partial_s = \partial_{x_2} + f_t'\partial_{y_1}$ (since the parametrisation in the $s$-direction moves in $(x_2,y_1)$). So the two tangent vectors to $\tilde{K}_t$ are $e_1 = \partial_{x_1}$ and $\partial_{x_2} + f_t'(s)\partial_{y_1}$.

Computing: $\omega(\partial_{x_1}, \partial_{x_2}) = 0$ (since $dx_1 \wedge dy_1$ gives $0$ on $(\partial_{x_1}, \partial_{x_2})$ and $dx_2 \wedge dy_2$ gives $0$). And $\omega(\partial_{x_1}, \partial_{y_1}) = dx_1 \wedge dy_1(\partial_{x_1},\partial_{y_1}) = 1$. Therefore
\[
  \omega(e_1, \partial_s) = 0 + f_t'(s) \cdot 1 = f_t'(s).
\]
This is nonzero in general, so the surface as parametrised above is \emph{not} Lagrangian. We need to correct the construction.

\medskip
\noindent\textbf{Corrected local construction.} The key is that in the standard hinge normal form, the edge $e$ lies along $e_1 = \partial_{x_1}$, $F_+$ lies in $\Pi_+ = \Span\{e_1, e_2\}$ (which has $y_1=y_2=0$), and $F_-$ lies in $\Pi_- = \Span\{e_1, e_4\} = \{x_2 = y_1 = 0\}$ (since both are Lagrangian and share $\ell = \Span\{e_1\}$, and the second half-plane is in the orthogonal complement in $\ell^\omega$). In coordinates $(x_1,x_2,y_1,y_2)$:
\[
  \Pi_+ = \{y_1 = y_2 = 0\}, \quad \Pi_- = \{x_2 = y_1 = 0\}.
\]
One checks: $\omega|_{\Pi_+} = (dx_1\wedge dy_1 + dx_2\wedge dy_2)|_{\{y_1=y_2=0\}} = 0$ (since $dy_1|_{\Pi_+} = dy_2|_{\Pi_+} = 0$). And $\omega|_{\Pi_-} = (dx_1\wedge dy_1 + dx_2\wedge dy_2)|_{\{x_2=y_1=0\}} = 0$ (since $dy_1|_{\Pi_-} = 0$ and $dx_2|_{\Pi_-} = 0$). So both are Lagrangian. Their intersection is $\{y_1=y_2=x_2=0\} = \Span\{e_1\}$. Good.

The local picture near $p$ in the $(x_2, y_2)$-plane (after projecting away $x_1$ and $y_1$) is: $F_+$ projects to $\{(x_2, 0) : x_2 \geq 0\}$ and $F_-$ projects to $\{(0, y_2) : y_2 \geq 0\}$. We wish to smooth the corner in the $(x_2, y_2)$-plane by a smooth curve, and then extend it over $x_1$.

Define, for $t > 0$ small, the smooth curve $\gamma_t$ in the $(x_2,y_2)$-half-plane ($x_2,y_2 \geq 0$) by
\begin{equation}
\label{eq:gamma_t}
  \gamma_t = \{(x_2, y_2) : x_2^2 + y_2^2 + t^2 = (x_2+y_2+t)^2 / 4 \text{... }\}.
\end{equation}
Actually, the simplest smoothing of two half-axes in $\mathbb{R}^2$ is the hyperbola $x_2 y_2 = t^2$ (for $x_2, y_2 > 0$), which is smooth and converges to the union of the axes as $t\to 0$. Observe that along $\gamma_t = \{x_2 y_2 = t^2, x_2, y_2 > 0\}$, the surface
\begin{equation}
\label{eq:Kt_local}
  \tilde{K}_t = \{(x_1, x_2, 0, y_2) : x_1 \in \mathbb{R},\, x_2 y_2 = t^2,\, x_2, y_2 > 0\}
\end{equation}
is an embedded surface lying in $\{y_1=0\}$. Its tangent vectors at $(x_1, x_2, 0, y_2)$ are $\partial_{x_1}$ and $\partial_s$ where $s$ parametrises $\gamma_t$.

On $\gamma_t$, differentiating $x_2 y_2 = t^2$ gives $y_2\, dx_2 + x_2\, dy_2 = 0$, so a tangent vector to $\gamma_t$ is $v = x_2 \partial_{x_2} - y_2 \partial_{y_2}$ (up to sign, writing $dy_2/dx_2 = -y_2/x_2$). Therefore the tangent plane to $\tilde{K}_t$ is $\Span\{\partial_{x_1}, x_2\partial_{x_2} - y_2\partial_{y_2}\}$. Computing:
\[
  \omega(\partial_{x_1}, x_2\partial_{x_2} - y_2\partial_{y_2})
  = x_2 \,\omega(\partial_{x_1},\partial_{x_2}) - y_2\,\omega(\partial_{x_1},\partial_{y_2})
  = x_2 \cdot 0 - y_2 \cdot 0 = 0.
\]
(We used $\omega(\partial_{x_1},\partial_{x_2}) = 0$ and $\omega(\partial_{x_1},\partial_{y_2}) = (dx_1\wedge dy_1 + dx_2\wedge dy_2)(\partial_{x_1},\partial_{y_2}) = 0 + 0 = 0$.) Hence $\tilde{K}_t$ is indeed Lagrangian. As $t\to 0$, $\tilde{K}_t$ converges to $\Pi_+^+ \cup \Pi_-^+$ (the union of the two positive half-planes), matching the original polyhedral surface $K$.

To promote this to a Hamiltonian isotopy, we note that $\tilde{K}_t$ are all exact Lagrangians in the following sense. Let $\lambda = \frac{1}{2}(x_1 dy_1 + x_2 dy_2 - y_1 dx_1 - y_2 dx_2)$ be the Liouville form (with $d\lambda = \omega$). On $\tilde{K}_t \subset \{y_1=0\}$, we compute $\lambda|_{\tilde{K}_t} = \frac{1}{2}(x_1 dy_1 - 0 + x_2 dy_2 - y_2 dx_2)|_{\tilde{K}_t}$. Since $y_1=0$ is constant on $\tilde{K}_t$, $dy_1|_{\tilde{K}_t}=0$. Also $x_2 dy_2 - y_2 dx_2 = x_2 dy_2 - y_2 dx_2$. On $\gamma_t$: $x_2 dy_2 - y_2 dx_2 = x_2(-y_2/x_2\, dx_2) - y_2 dx_2 = -2y_2 dx_2$, and $x_1 dy_1|_{\tilde{K}_t}=0$. So $\lambda|_{\tilde{K}_t} = \frac{1}{2}(-2y_2 dx_2) = -y_2 dx_2$. Since $y_2 = t^2/x_2$, we get $\lambda|_{\tilde{K}_t} = -(t^2/x_2)dx_2 = -d(t^2 \log x_2)$, an exact 1-form. Hence $\lambda|_{\tilde{K}_t}$ is exact, so each $\tilde{K}_t$ is an exact Lagrangian.

By the Lagrangian isotopy--Hamiltonian isotopy correspondence (Banyaga's theorem, or directly: since $\mathbb{R}^4$ is simply connected and $\tilde{K}_t$ is a smooth isotopy of exact Lagrangians for $t > 0$, by Weinstein's neighborhood theorem and Moser's trick, the isotopy $\{\tilde{K}_t\}$ for $t \in (0,1]$ is realised by a Hamiltonian isotopy of $\mathbb{R}^4$ supported near the edge), the family $\tilde{K}_t$ is a Hamiltonian isotopy of smooth Lagrangians for $t > 0$, and with $\tilde{K}_0 = K$ locally near $e$, this is the required Lagrangian smoothing near $e$.

To make this global (and supported in $B(p,2\varepsilon)$): multiply the smoothing parameter $t$ by a smooth cutoff function $\chi(x_1)$ supported in $(-\varepsilon, \varepsilon)$ with $\chi \equiv 1$ on $(-\varepsilon/2, \varepsilon/2)$. The resulting surface is Lagrangian (the Lagrangian condition is local in $x_1$, and at each fixed $x_1$ the cross-section is a Lagrangian curve in the $(x_2,y_2)$-plane), smooth for $t > 0$, and agrees with $K$ outside $B(p, 2\varepsilon)$.
\end{proof}

\section*{Step 2: Smoothing of a vertex}

\begin{lemma}[Vertex smoothing]
\label{lem:vertex}
Let $v$ be a vertex of $K$ at which exactly four faces $F_1, F_2, F_3, F_4$ (cyclically ordered) meet. Let $\varepsilon > 0$ be small enough that $B(v,\varepsilon)$ contains no other vertex. Then there exists a smooth family $\{K_t\}_{t \in [0,1]}$ such that:
\begin{enumerate}
  \item[(a)] $K_0 = K$,
  \item[(b)] for each $t \in (0,1]$, $K_t \cap B(v, \varepsilon)$ is a smooth embedded Lagrangian,
  \item[(c)] $K_t = K$ outside $B(v,2\varepsilon)$,
  \item[(d)] the isotopy $\{K_t\}_{t \in (0,1]}$ is Hamiltonian.
\end{enumerate}
\end{lemma}

\begin{proof}
By Lemma \ref{lem:vertex_model}, after a linear symplectomorphism and translation taking $v$ to the origin, the four Lagrangian tangent planes at $v$ are in the standard form \eqref{eq:standard_vertex}:
\[
  \Pi_1 = \Span\{e_1,e_2\},\quad \Pi_2 = \Span\{e_2,e_3\},\quad
  \Pi_3 = \Span\{e_3,e_4\},\quad \Pi_4 = \Span\{e_4,e_1\},
\]
and the four adjacent edges lie along $\ell_{12} = \Span\{e_2\}$, $\ell_{23}=\Span\{e_3\}$, $\ell_{34}=\Span\{e_4\}$, $\ell_{41}=\Span\{e_1\}$.

We construct an explicit smooth Lagrangian surface that replaces the cone $C = \Pi_1^+ \cup \Pi_2^+ \cup \Pi_3^+ \cup \Pi_4^+$ (four quarter-plane sectors) near the origin.

\medskip
\noindent\textbf{Product structure.} Observe that the standard symplectic form $\omega = dx_1\wedge dy_1 + dx_2\wedge dy_2$ makes $\mathbb{R}^4 = \mathbb{R}^2_{x_1,y_1} \times \mathbb{R}^2_{x_2,y_2}$ a product of two symplectic planes. In this product:
\[
  \Pi_1 = \mathbb{R}_{x_1} \times \mathbb{R}_{x_2},\quad
  \Pi_2 = \mathbb{R}_{x_2} \times \mathbb{R}_{y_1},\quad
  \Pi_3 = \mathbb{R}_{y_1} \times \mathbb{R}_{y_2},\quad
  \Pi_4 = \mathbb{R}_{y_2} \times \mathbb{R}_{x_1},
\]
where we use the natural identification: $e_1=\partial_{x_1}$, $e_2=\partial_{x_2}$, $e_3=\partial_{y_1}$, $e_4=\partial_{y_2}$.

Let $\gamma_1, \gamma_2$ be smooth embedded curves in $\mathbb{R}^2$. Consider the product surface $\Sigma = \gamma_1 \times \gamma_2 \subset \mathbb{R}^4$. This is Lagrangian if and only if the product of the curves is Lagrangian for the product symplectic structure. For $\Sigma = \gamma_1 \times \gamma_2$ with tangent vectors $(\dot\gamma_1(s_1), 0)$ and $(0, \dot\gamma_2(s_2))$, we compute:
\[
  \omega\bigl((\dot\gamma_1, 0), (0, \dot\gamma_2)\bigr)
  = \omega_1(\dot\gamma_1, 0) + \omega_2(0, \dot\gamma_2) = 0,
\]
where $\omega_1 = dx_1 \wedge dy_1$ and $\omega_2 = dx_2 \wedge dy_2$. Since the two tangent vectors are in orthogonal factors, $\omega(\cdot,\cdot) = 0$ identically, and $\Sigma$ is automatically isotropic. Since $\dim\Sigma = 2 = \frac{1}{2}\dim\mathbb{R}^4$, $\Sigma$ is Lagrangian.

\medskip
\noindent\textbf{Choosing the profile curves.} We choose $\gamma_1$ and $\gamma_2$ so that $\gamma_1 \times \gamma_2$ approximates the cone $C$ near the origin and is smooth.

Consider the figure-eight curve $\gamma_1$ in $\mathbb{R}^2_{x_1,y_1}$ and similarly $\gamma_2$ in $\mathbb{R}^2_{x_2,y_2}$. However, we need to be more careful: the four planes $\Pi_i$ are not a simple product of curves; they are a cyclic arrangement. Let us see which cross-sections they define.

The four face-planes are: $\Pi_1 = \{y_1=y_2=0\}$, $\Pi_2 = \{x_1=y_2=0\}$, $\Pi_3 = \{x_1=x_2=0\}$, $\Pi_4 = \{x_2=y_1=0\}$.

Consider the intersection of $C = \Pi_1^+ \cup \Pi_2^+ \cup \Pi_3^+ \cup \Pi_4^+$ with the unit sphere $S^3$. This is a closed curve (the \emph{link} of the vertex), which is a piecewise-great-circle curve. More precisely, on $S^3 = \{x_1^2+x_2^2+y_1^2+y_2^2=1\}$, the curve $C \cap S^3$ has four arcs:
\begin{align*}
  C_1 &= \Pi_1 \cap S^3 \cap \{x_1\geq 0, x_2\geq 0\} = \{(x_1,x_2,0,0): x_1^2+x_2^2=1,\, x_1,x_2\geq 0\},\\
  C_2 &= \Pi_2 \cap S^3 \cap \{x_2\geq 0, y_1\geq 0\} = \{(0,x_2,y_1,0): x_2^2+y_1^2=1,\, x_2,y_1\geq 0\},\\
  C_3 &= \Pi_3 \cap S^3 \cap \{y_1\geq 0, y_2\geq 0\} = \{(0,0,y_1,y_2): y_1^2+y_2^2=1,\, y_1,y_2\geq 0\},\\
  C_4 &= \Pi_4 \cap S^3 \cap \{y_2\geq 0, x_1\geq 0\} = \{(x_1,0,0,y_2): x_1^2+y_2^2=1,\, x_1,y_2\geq 0\},
\end{align*}
forming a closed piecewise-smooth curve $\Lambda = C_1 \cup C_2 \cup C_3 \cup C_4$ in $S^3$.

We must verify that $\Lambda$ is a Legendrian unknot in $(S^3, \xi_{std})$, the standard contact structure $\xi_{std} = \ker(\alpha)$ where $\alpha = \frac{1}{2}(x_1 dy_1 - y_1 dx_1 + x_2 dy_2 - y_2 dx_2)|_{S^3}$.

\medskip
\noindent\textbf{$\Lambda$ is Legendrian.} The contact form $\alpha$ restricts to the Liouville $1$-form on $S^3$. A curve is Legendrian if its tangent vectors lie in $\ker\alpha$.

On arc $C_1 = \{(x_1,x_2,0,0): x_1^2+x_2^2=1\}$: tangent vector is $\dot C_1 = (-x_2, x_1, 0, 0)$ (parametrising by angle). Then $\alpha(\dot C_1) = \frac{1}{2}(x_1 \cdot 0 - 0 \cdot (-x_2) + x_2 \cdot 0 - 0 \cdot x_1) = 0$. So $C_1$ is Legendrian.

On arc $C_2 = \{(0,x_2,y_1,0): x_2^2+y_1^2=1\}$: tangent vector is $(-y_1, x_2, 0, 0)$ reparametrised... The parametrisation by angle $\varphi$: $(x_2,y_1) = (\cos\varphi, \sin\varphi)$, so tangent vector in $\mathbb{R}^4$ is $(0,-\sin\varphi, \cos\varphi, 0)$. Then $\alpha = \frac{1}{2}(x_1 dy_1 - y_1 dx_1 + x_2 dy_2 - y_2 dx_2)$. Here $x_1=0, y_2=0$, so $\alpha = \frac{1}{2}(-y_1 dx_1 + x_2 dy_2)|_{C_2}$. Since on $C_2$ both $x_1=0$ and $y_2=0$ are constant, $dx_1=dy_2=0$ along $C_2$, so $\alpha|_{C_2}=0$ and all tangent vectors lie in $\ker\alpha|_{C_2}$. Hence $C_2$ is Legendrian.

By analogous calculations (permuting coordinates cyclically), $C_3$ and $C_4$ are also Legendrian (each arc lies in a $2$-plane through the origin on which $\alpha$ vanishes identically). Hence $\Lambda = C_1 \cup C_2 \cup C_3 \cup C_4$ is a piecewise-Legendrian curve.

\medskip
\noindent\textbf{$\Lambda$ is topologically unknotted.} As a closed curve in $S^3$, $\Lambda$ consists of four quarter-circle arcs joined end to end, forming a closed loop. This is isotopic (as a topological curve) to the standard unknot: $\Lambda$ bounds an obvious disk, namely the surface $\{(x_1,x_2,0,0): x_1^2+x_2^2 \leq 1, x_1,x_2 \geq 0\}$ union the analogous quarter-disks in $\Pi_2,\Pi_3,\Pi_4$... In fact $\Lambda$ is clearly contractible in $S^3$ (it is a simple closed curve on the boundary of a topological ball in $S^3$). More concretely, $\Lambda$ can be continuously deformed to a single great circle by a homotopy within $S^3$, so it is unknotted.

The Thurston--Bennequin invariant of $\Lambda$ as a Legendrian knot can be computed from its front projection or by counting linking numbers. For the standard model, a direct computation gives $\text{tb}(\Lambda) = -1$: since $\Lambda$ is a Legendrian unknot with $\text{tb}=-1$, it bounds a Lagrangian disk by the standard construction (the Lagrangian filling of a Legendrian unknot with $\text{tb}=-1$).

\medskip
\noindent\textbf{Explicit smooth Lagrangian.} We now construct the explicit smooth Lagrangian surface $\tilde{K}_t$ that replaces $C$ near the origin.

Using the product structure, define for $t > 0$:
\begin{equation}
\label{eq:product_surface}
  \tilde{K}_t = \{(x_1, x_2, y_1, y_2) \in \mathbb{R}^4 :
    (x_1^2 + y_1^2)(x_2^2 + y_2^2) = t^4, \;
    x_1 y_1 \geq 0, \; x_2 y_2 \geq 0\}.
\end{equation}
This is the product of two level sets: in the $(x_1,y_1)$-plane, consider the curve $\Gamma_1 = \{r_1 = t^2/r_2\}$ where $r_i = \sqrt{x_i^2+y_i^2}$, and similarly in the $(x_2,y_2)$-plane. More precisely, $\tilde{K}_t$ is parametrised by $(r_1, \theta_1, r_2, \theta_2)$ where $r_1 r_2 = t^2$, $\theta_i \in [0,\pi/2]$ (restricting to positive octants to match the face orientations), via
\[
  (x_1, y_1, x_2, y_2) = (r_1 \cos\theta_1, r_1\sin\theta_1, r_2\cos\theta_2, r_2\sin\theta_2).
\]
The surface $\tilde{K}_t$ is parametrised by $(r_1, \theta_1, \theta_2)$ with $r_2 = t^2/r_1$. As $t\to 0$, either $r_1\to 0$ (and then $r_2\to\infty$) or $r_2\to 0$, so the surface degenerates to the four Lagrangian planes.

However, for the local smoothing near the origin, we restrict to a compact region $\{r_1^2+r_2^2 \leq R^2\}$ and construct a compact surface. Instead, let us use a more explicit product construction.

Define smooth curves in $\mathbb{R}^2$:
\begin{equation}
\label{eq:curves}
  \gamma_t^{(1)} = \{(x_1, y_1) : x_1 y_1 = t^2,\; x_1,y_1 > 0\} \cup \{(x_1,y_1): x_1 y_1 = t^2,\; x_1,y_1 < 0\},
\end{equation}
i.e.\ the hyperbola $\{x_1 y_1 = t^2\}$ in both the first and third quadrants. Similarly let $\gamma_t^{(2)} = \{(x_2,y_2): x_2 y_2 = t^2,\; x_2,y_2>0\}$.

The product $\gamma_t^{(1)} \times \gamma_t^{(2)} \subset \mathbb{R}^4$ is a smooth Lagrangian surface for any $t > 0$, by the product Lagrangian argument above (since $\omega(\dot\gamma^{(1)},0) = 0$ and $\omega(0,\dot\gamma^{(2)}) = 0$ trivially, and $\omega(\dot\gamma^{(1)}, \dot\gamma^{(2)}) = \omega_1(\dot\gamma^{(1)},0) + \omega_2(0,\dot\gamma^{(2)}) = 0$). Wait, this is not the right pairing; let us restate: the tangent space of $\gamma_t^{(1)} \times \gamma_t^{(2)}$ at $(p,q)$ is $T_p\gamma^{(1)} \oplus T_q\gamma^{(2)} \subset \mathbb{R}^2 \oplus \mathbb{R}^2 = \mathbb{R}^4$. For vectors $v = (v_1,0) \in T_p\gamma^{(1)}\oplus 0$ and $w=(0,w_2) \in 0\oplus T_q\gamma^{(2)}$, $\omega(v,w) = \omega_1(v_1,0) + \omega_2(0,w_2)$; but $\omega_1(v_1, 0) = 0$ since $\omega_1$ is a $2$-form on $\mathbb{R}^2_{x_1,y_1}$ and $0\in\mathbb{R}^2$ is a zero vector; similarly $\omega_2(0,w_2)=0$. And for two vectors both in the first factor: $\omega((v_1,0),(u_1,0)) = \omega_1(v_1,u_1)$, which is $0$ if and only if $\gamma^{(1)}$ is isotropic in $(\mathbb{R}^2,\omega_1)$. But any $1$-manifold in a $2$-dimensional symplectic space is Lagrangian (isotropic of half-dimension), so $\omega_1|_{T\gamma^{(1)}} = 0$. Similarly $\omega_2|_{T\gamma^{(2)}} = 0$. Hence $\omega|_{T(\gamma^{(1)}\times\gamma^{(2)})} = 0$, confirming $\gamma_t^{(1)}\times\gamma_t^{(2)}$ is Lagrangian.

\medskip
\noindent\textbf{Convergence to the polyhedral surface.} The hyperbola $\{x_1 y_1 = t^2, x_1,y_1>0\}$ converges, as $t\to 0^+$, to the two coordinate half-axes $\{x_1>0,y_1=0\}\cup\{x_1=0,y_1>0\}$ in the Hausdorff sense on compact subsets (avoiding the origin). More precisely, for any compact $D \subset \mathbb{R}^2\setminus\{0\}$, the restriction of $\gamma_t^{(1)}$ to $D$ converges in $C^\infty$ topology to the half-axes as $t\to 0$.

Similarly $\gamma_t^{(2)}$ converges to $\{x_2>0,y_2=0\}\cup\{x_2=0,y_2>0\}$. The product $\gamma_t^{(1)}\times\gamma_t^{(2)}$ therefore converges on compact subsets of $\mathbb{R}^4\setminus C$ to four sheets which approach the four faces of $K$ near the origin: the four products of the half-axis branches.

More carefully: $\gamma_t^{(1)}$ has two branches $\gamma_t^{(1),+} = \{x_1y_1=t^2, x_1,y_1>0\}$ and $\gamma_t^{(1),-} = \{x_1y_1=t^2, x_1,y_1<0\}$. Similarly for $\gamma_t^{(2)}$. For the local model near $v$, we restrict to $\gamma_t^{(1),+}\times\gamma_t^{(2),+}$, which is a single connected smooth surface. This converges to the two half-axes from $\gamma^{(1),+}$ (the positive $x_1$ and positive $y_1$ half-axes) crossed with the two half-axes from $\gamma^{(2),+}$ (positive $x_2$ and positive $y_2$ half-axes), giving four half-planes matching $F_1\cap\ldots\cup F_4$.

To see this: $\gamma_t^{(1),+}\times\gamma_t^{(2),+}$ consists of points $(x_1,y_1,x_2,y_2)$ with $x_1y_1=t^2>0$ and $x_2y_2=t^2>0$. The four branches of the limit (as $t\to 0$ on compact sets away from coordinate axes) approach:
\begin{itemize}
  \item $\{x_1\to\infty, y_1\to 0, x_2\to\infty, y_2\to 0\}$: the face $\Pi_1^+$,
  \item $\{x_1\to\infty, y_1\to 0, y_2\to\infty, x_2\to 0\}$: the face $\Pi_4^+$,
  \item $\{y_1\to\infty, x_1\to 0, x_2\to\infty, y_2\to 0\}$: the face $\Pi_2^+$,
  \item $\{y_1\to\infty, x_1\to 0, y_2\to\infty, x_2\to 0\}$: the face $\Pi_3^+$,
\end{itemize}
which are exactly the four faces $\Pi_1^+, \Pi_2^+, \Pi_3^+, \Pi_4^+$ of the cone $C$.

\medskip
\noindent\textbf{Compactification and cutoff.} To make this local (supported near $v$), introduce a smooth bump: let $\chi\colon\mathbb{R}^4\to[0,1]$ be supported in $B(0,2\varepsilon)$ with $\chi\equiv 1$ on $B(0,\varepsilon)$. Define the smoothing parameter $t = t(\varepsilon) > 0$ small (to be chosen). The actual smoothed surface $K_t^{(v)}$ near $v$ is obtained by deforming only within $B(0,\varepsilon)$: we replace $C\cap B(0,\varepsilon)$ with the corresponding portion of $\gamma_t^{(1),+}\times\gamma_t^{(2),+}$, and match back to $K$ near $\partial B(0,\varepsilon)$ via a smooth transition. This can be done because for $t$ small relative to $\varepsilon$, the surface $\gamma_t^{(1),+}\times\gamma_t^{(2),+}$ is uniformly close to the cone $C$ on $\partial B(0,\varepsilon)$, and a standard $C^\infty$ interpolation preserves the Lagrangian condition (by the Lagrangian stability theorem: nearby surfaces can be made Lagrangian by a small Hamiltonian perturbation, since $\omega$ restricts to zero on each face of $C$).

\medskip
\noindent\textbf{Hamiltonian realisation.} We verify the isotopy is Hamiltonian. The family $\{K_t^{(v)}\}_{t\in(0,1]}$ is a smooth isotopy of compact Lagrangian surfaces in $\mathbb{R}^4$. Since $\mathbb{R}^4$ is simply connected, the flux of the isotopy vanishes (the flux homomorphism factors through $H^1(K_t;\mathbb{R}) \to H^1(\mathbb{R}^4;\mathbb{R}) = 0$). By the Banyaga theorem (see e.g.\ Salamon \cite{salamon}), a Lagrangian isotopy in a symplectic manifold is Hamiltonian if and only if its flux vanishes. Hence $\{K_t^{(v)}\}$ is Hamiltonian.

More concretely: on the product surface $\gamma_t^{(1),+}\times\gamma_t^{(2),+}$, the Liouville form $\lambda = \frac{1}{2}\sum_i(x_i dy_i - y_i dx_i)$ restricts to
\[
  \lambda|_{\gamma_t^{(1),+}\times\gamma_t^{(2),+}}
  = \frac{1}{2}(x_1 dy_1 - y_1 dx_1)|_{\gamma_t^{(1),+}}
  + \frac{1}{2}(x_2 dy_2 - y_2 dx_2)|_{\gamma_t^{(2),+}}.
\]
On $\gamma_t^{(1),+} = \{x_1y_1=t^2\}$: differentiating gives $y_1 dx_1 + x_1 dy_1 = 0$, so $dy_1 = -\frac{y_1}{x_1}dx_1$. Then
$x_1 dy_1 - y_1 dx_1 = x_1(-\frac{y_1}{x_1}dx_1) - y_1 dx_1 = -2y_1 dx_1 = -2\frac{t^2}{x_1}dx_1 = -d(2t^2\log x_1)$.
So $\lambda|_{\gamma_t^{(1),+}} = -d(t^2\log x_1)$. Similarly $\lambda|_{\gamma_t^{(2),+}} = -d(t^2\log x_2)$. Hence $\lambda|_{\gamma_t^{(1),+}\times\gamma_t^{(2),+}} = -d(t^2(\log x_1+\log x_2)) = -d(t^2\log(x_1 x_2))$, which is exact. So each $\gamma_t^{(1),+}\times\gamma_t^{(2),+}$ is an exact Lagrangian submanifold of $\mathbb{R}^4$, confirming the flux is zero and the isotopy is Hamiltonian.
\end{proof}

\section*{Step 3: Global smoothing}

\begin{proof}[Proof of Theorem \ref{thm:main}]
Let $K$ be a polyhedral Lagrangian surface with exactly $4$ faces at every vertex. Let $V = \{v_1,\ldots,v_m\}$ be the set of vertices and $E = \{e_1,\ldots,e_k\}$ be the set of edges. Since $K$ is a finite polyhedral complex, $V$ and $E$ are finite sets.

\medskip
\noindent\textbf{Step A: Choose disjoint neighborhoods.} Since $V$ is finite, choose $\varepsilon > 0$ small enough so that the balls $\{B(v_i, 2\varepsilon)\}_{i=1}^m$ are pairwise disjoint, and each $B(v_i,2\varepsilon)$ contains no edge-interior point from any edge not adjacent to $v_i$.

For each edge $e_j$, the edge $e_j$ minus the two balls around its endpoints, $e_j \setminus \bigcup_i B(v_i,\varepsilon)$, is a compact interval in the interior of the edge. For each interior point $p$ of this reduced edge, choose a ball $B(p, \varepsilon_p)$ small enough that it is disjoint from all vertex balls and from all other edges. By compactness, finitely many such balls cover $e_j \setminus \bigcup_i B(v_i,\varepsilon)$; call them $B_{j,1},\ldots,B_{j,n_j}$.

\medskip
\noindent\textbf{Step B: Local smoothings.}
\begin{enumerate}
  \item For each vertex $v_i$, apply Lemma \ref{lem:vertex} with radius $\varepsilon$ to obtain a Hamiltonian isotopy $\{K_t^{(v_i)}\}_{t\in[0,1]}$ supported in $B(v_i,2\varepsilon)$, smoothing $K$ near $v_i$.
  \item For each edge ball $B_{j,l}$, apply Lemma \ref{lem:edge} to obtain a Hamiltonian isotopy $\{K_t^{(j,l)}\}_{t\in[0,1]}$ supported in $2B_{j,l}$, smoothing $K$ near the corresponding edge point.
\end{enumerate}

\medskip
\noindent\textbf{Step C: Combining the local smoothings.}
Since all the support balls are pairwise disjoint (by choice of $\varepsilon$ and the finite cover), we can compose all the local Hamiltonian isotopies simultaneously. More precisely: since the Hamiltonian functions for vertex and edge smoothings are supported in disjoint open sets, their flows commute and we may define the combined Hamiltonian $H = \sum_i H^{(v_i)} + \sum_{j,l} H^{(j,l)}$, where $H^{(v_i)}$ (resp.\ $H^{(j,l)}$) are the Hamiltonian functions generating the vertex (resp.\ edge) smoothings. The time-$t$ flow of $H$ is the product of the individual flows, which are the identity outside their respective support balls.

Setting $K_t = \phi_t^H(K)$ for $t\in(0,1]$ and $K_0 = K$, we obtain:
\begin{itemize}
  \item For $t\in(0,1]$, $K_t$ is a smooth embedded Lagrangian: near each vertex $v_i$, the vertex smoothing gives a smooth Lagrangian piece; along each edge (away from vertices), the edge smoothing gives a smooth Lagrangian piece; on the faces (away from edges), $K_t = K$ and the faces are smooth flat Lagrangians. These pieces fit together smoothly because the support balls are disjoint.
  \item As $t\to 0$, each $K_t^{(v_i)} \to K$ near $v_i$ and each $K_t^{(j,l)} \to K$ near the edge points, so $K_t \to K$ uniformly, giving the topological isotopy with $K_0 = K$.
  \item The isotopy $\{K_t\}_{t\in(0,1]}$ is Hamiltonian, being generated by $H$.
\end{itemize}
Thus $\{K_t\}_{t\in[0,1]}$ is the required Lagrangian smoothing.
\end{proof}

\begin{remark}
The key role of the \emph{$4$-valency hypothesis} is in Lemma \ref{lem:vertex_model}: it allows the four tangent planes at a vertex to be simultaneously normalised to the standard product form \eqref{eq:standard_vertex}. With this product structure, the product surface $\gamma_t^{(1),+}\times\gamma_t^{(2),+}$ is automatically Lagrangian without any additional constraint. For vertices of valency $\neq 4$, the configuration of tangent planes need not admit such a product structure, and additional obstructions (encoded in the Maslov class and the Thurston--Bennequin number of the Legendrian link) may prevent smoothing.

More precisely, at a $k$-valent vertex, the Legendrian link $\Lambda = C\cap S^3$ is a closed curve with $k$ Legendrian arcs. For $k=4$, as shown above, $\Lambda$ is a Legendrian unknot with $\text{tb}(\Lambda)=-1$, which is the maximal Thurston--Bennequin invariant for the unknot. This is the condition for a Legendrian unknot to bound a Lagrangian disk, and hence for the cone singularity to be Lagrangian fillable (smoothable). For $k\neq 4$, the Thurston--Bennequin invariant may be lower (or the knot type different), blocking such a filling.
\end{remark}

\begin{thebibliography}{99}

\bibitem{salamon}
D.~Salamon,
\emph{Lectures on Floer homology},
in: Symplectic geometry and topology (Park City, UT, 1997),
IAS/Park City Math. Ser.\ \textbf{7}, Amer.\ Math.\ Soc., Providence, RI, 1999, pp.~143--229.

\bibitem{weinstein}
A.~Weinstein,
\emph{Symplectic manifolds and their Lagrangian submanifolds},
Adv.\ Math.\ \textbf{6} (1971), 329--346.

\bibitem{polterovich}
L.~Polterovich,
\emph{The surgery of Lagrange submanifolds},
Geom.\ Funct.\ Anal.\ \textbf{1} (1991), no.~2, 198--210.

\bibitem{banyaga}
A.~Banyaga,
\emph{Sur la structure du groupe des diff\'{e}omorphismes qui pr\'{e}servent une forme symplectique},
Comment.\ Math.\ Helv.\ \textbf{53} (1978), no.~2, 174--227.

\bibitem{arnold}
V.~I.~Arnold,
\emph{Mathematical Methods of Classical Mechanics},
2nd ed., Graduate Texts in Mathematics, Vol.~60,
Springer-Verlag, New York, 1989.

\bibitem{mcduff_salamon}
D.~McDuff and D.~Salamon,
\emph{Introduction to Symplectic Topology},
3rd ed., Oxford University Press, 2017.

\end{thebibliography}

\end{document}
